\ncdmodultestmthree{2}{Sundhed, Fritid og Digitalisering}{全真模考卷二：数字排毒、酒店投诉与健康习惯}{一周无社交媒体挑战 / 合买纪念礼物与膝伤休跑聊天 / 社区美食社团 / 三人谈手机与屏幕时间 / 酒店投诉信与夜间提升课程回信 / 健康习惯与数字生活思维导图及运动休闲四图互动}{复合从句连词与词序（fordi/før/at 否定与状语前置）· 跨回合指代与语气衔接 · 因果/转折句段内前瞻后顾 · 被动语态（-s vs blive）· 介词 + 形式宾语 det + 从句 · 书面复合介词（på baggrund af / til trods for）}

% =====================================================================
% DEL 1: LÆSNING
% =====================================================================
\ncdsubsec{Del 1: Læsning（阅读理解｜50 minutter · 24 point）}

\begin{tipbox}{考场规则与配速建议（Læsning｜15--24 point = Bestået）}
\small
本部分共五道模块题，考场将 \textbf{Opgave 1, 2A, 2B, 3 og 4} 一次性合订发放。限时 \textbf{50 分钟}，全程\textbf{禁用任何词典、笔记或电子设备}。每题 1 分，满分 \textbf{24 分}，及格线为 \textbf{15 分}（0--14 分 = \textit{Ikke bestået}）。建议配速：Opgave 1（8 分钟）$\rightarrow$ Opgave 2A \& 2B（10 分钟）$\rightarrow$ Opgave 3（14 分钟）$\rightarrow$ Opgave 4（13 分钟），最后留 5 分钟复查。
\end{tipbox}

% ---------------------------------------------------------------------
\subsubsection*{Opgave 1 (6 point)：En uge uden sociale medier}

\textbf{Instruktion:} Læs teksten \textit{En uge uden sociale medier}. Skriv de ord (1--6), der mangler. Du skal kun bruge hvert ord én gang. Der er seks ord, du ikke skal bruge. Se eksemplet (0).

\vspace{0.35em}
\textbf{\color{NCDNavy}En uge uden sociale medier}

\vspace{0.2em}
Mikkel er niogtyve år og arbejder som grafisk designer i Roskilde. I mange år brugte han næsten tre timer hver dag på sociale medier, og (0)\ \underline{\textbf{det}} begyndte langsomt at irritere ham. Når vækkeuret ringede om morgenen, tjekkede han straks sine beskeder, (1)\ \rule{2.0cm}{0.4pt} han overhovedet stod ud af sengen. På arbejdet mistede han ofte fokus, (2)\ \rule{2.0cm}{0.4pt} skærmen hele tiden lyste med nye notifikationer og videoklip. Derfor besluttede han sidste måned at lave et eksperiment: Han ville slette alle sociale apps fra telefonen i syv dage for at mærke forskellen.

De første to døgn var overraskende hårde. Når Mikkel sad i toget eller ventede på kaffen i kantinen, søgte hånden næsten automatisk ned i lommen efter mobilen. Efter tre dage blev alting (3)\ \rule{2.0cm}{0.4pt} meget roligere. Uden den konstante strøm af billeder sov han bedre om natten og fik pludselig tid (4)\ \rule{2.0cm}{0.4pt} at læse en roman, som havde ligget på natbordet i et halvt år. Han mærkede også tydeligt, (5)\ \rule{2.0cm}{0.4pt} han talte mere med sine kolleger i pauserne i stedet for at stirre ned i skærmen. Da ugen endelig var slut, hentede Mikkel kun to apps tilbage og slog alle notifikationer fra. Nu bruger han højst tyve minutter om dagen på nettet, og det har (6)\ \rule{2.0cm}{0.4pt} ham mere energi i hverdagen.

\vspace{0.35em}
\begin{grammatikbox}{Ordboks（选词框 · 共 13 词：例词 (0) 已用 + 6 答案 + 6 多余干扰词）}
\centering
\large\bfseries\color{NCDNavy}
at \qquad blev \qquad {\color{NCDGray}\textit{det}\textsuperscript{(0)}} \qquad dog \qquad eller \qquad fordi \qquad fra \\[3pt]
før \qquad givet \qquad med \qquad men \qquad som \qquad til
\end{grammatikbox}

% ---------------------------------------------------------------------
\subsubsection*{Opgave 2A (NY · 3 point)：Lukas skriver til sin veninde Emma}

\textbf{Instruktion:}\\
\textbf{Situation:} Lukas skriver til sin veninde Emma.
\begin{itemize}
  \item Læs Lukas' og Emmas chat. Der er 5 beskeder fra Lukas.
  \item Find de svar fra Emma, der passer (A--F).
  \item Skriv bogstavet på det svar, der passer.
\end{itemize}
\hspace*{1.5em}Der er tre svar, du ikke skal bruge.\\
\hspace*{1.5em}Se eksemplet (0).

\vspace{0.25em}
\small
\renewcommand{\arraystretch}{1.35}
\begin{longtable}{@{} >{\raggedright\arraybackslash}p{8.1cm} @{\hspace{0.5cm}} >{\raggedright\arraybackslash}p{8.1cm} @{}}
\toprule
\textbf{\color{NCDNavy}Chat mellem Lukas og Emma} & \textbf{\color{NCDNavy}Svar fra Emma (Z \& A--F)} \\
\midrule
\endfirsthead
\toprule
\textbf{\color{NCDNavy}Chat mellem Lukas og Emma（续表）} & \textbf{\color{NCDNavy}Svar fra Emma（续表）} \\
\midrule
\endhead
\bottomrule
\endfoot
\badge{NCDNavy}{Lukas 1}\ Hej Emma! Signe og Kasper har jo fem års jubilæum næste lørdag. Har du fundet på en god gave fra os alle sammen endnu?\par\vspace{4pt}
\textbf{0. Emma:}\ \framebox[0.85cm]{\textbf{Z}} \quad \textcolor{NCDGray}{\small (Eksempel)} &
\badge{NCDGray}{Z}\ \textcolor{NCDGray}{\textit{Hej! Ikke helt endnu, men jeg tænkte på et gavekort til en god restaurant. Hvad siger du til det?}} \\
\midrule
\badge{NCDNavy}{Lukas 2}\ Ja, det er en super idé! De elsker jo begge to italiensk mad. Hvor meget synes du, at hver person skal betale til gaven?\par\vspace{4pt}
\textbf{1. Emma:}\ \framebox[0.85cm]{\rule{0pt}{0.26cm}} &
\badge{NCDNavy}{A}\ Nej, de aner ingenting! Kaspers bror lukker os stille ind kl. 14.30, mens de er ude at handle.\par\vspace{6pt}
\badge{NCDNavy}{B}\ Hvad med 100 kroner hver? Så bliver det ikke for dyrt for nogen, og vi skal jo også have noget lækkert til kaffen. \\
\midrule
\badge{NCDNavy}{Lukas 3}\ Det lyder helt fair, så bliver det i alt 800 kroner, hvis vi er otte personer. Skal jeg også bestille kagen hos bageren?\par\vspace{4pt}
\textbf{2. Emma:}\ \framebox[0.85cm]{\rule{0pt}{0.26cm}} &
\badge{NCDNavy}{C}\ Nej, vi behøver slet ikke nogen kage, for vi har allerede spist for meget sukker i denne uge.\par\vspace{6pt}
\badge{NCDNavy}{D}\ Det er ikke nødvendigt at give 100 kroner hver, for jeg betaler bare hele gaven og kagen alene. \\
\midrule
\badge{NCDNavy}{Lukas 4}\ Okay, så ringer jeg til bageriet i morgen formiddag og bestiller en chokoladekage uden nødder. Hvad med overraskelsen -- ved de, at vi kommer forbi lørdag eftermiddag?\par\vspace{4pt}
\textbf{3. Emma:}\ \framebox[0.85cm]{\rule{0pt}{0.26cm}} &
\badge{NCDNavy}{E}\ Ja, det må du meget gerne gøre, for du bor lige ved siden af bageriet. Kan du bestille en chokoladekage uden nødder?\par\vspace{6pt}
\badge{NCDNavy}{F}\ Jo, de ved godt, at vi kommer kl. 15, for jeg har allerede sendt dem en invitation. \\
\midrule
\badge{NCDNavy}{Lukas 5}\ Perfekt, så holder jeg det selvfølgelig helt hemmeligt, og jeg skriver lige til de andre gæster om mødetidspunktet. Vi ses på lørdag kl. 14.30! &
\textcolor{NCDGray}{\small （第 5 条为 Lukas 的收尾回复，无需补全）} \\
\end{longtable}
\normalsize

% ---------------------------------------------------------------------
\subsubsection*{Opgave 2B (NY · 3 point)：Sofie skriver til sin løbetræner Peter}

\textbf{Instruktion:}\\
\textbf{Situation:} Sofie skriver til sin løbetræner Peter.
\begin{itemize}
  \item Læs Sofies og Peters chat. Der er 5 beskeder fra Sofie.
  \item Find de svar fra Peter, der passer (A--F).
  \item Skriv bogstavet på det svar, der passer.
\end{itemize}
\hspace*{1.5em}Der er tre svar, du ikke skal bruge.\\
\hspace*{1.5em}Se eksemplet (0).

\vspace{0.25em}
\small
\renewcommand{\arraystretch}{1.35}
\begin{longtable}{@{} >{\raggedright\arraybackslash}p{8.1cm} @{\hspace{0.5cm}} >{\raggedright\arraybackslash}p{8.1cm} @{}}
\toprule
\textbf{\color{NCDNavy}Chat mellem Sofie og Peter} & \textbf{\color{NCDNavy}Svar fra Peter (Z \& A--F)} \\
\midrule
\endfirsthead
\toprule
\textbf{\color{NCDNavy}Chat mellem Sofie og Peter（续表）} & \textbf{\color{NCDNavy}Svar fra Peter（续表）} \\
\midrule
\endhead
\bottomrule
\endfoot
\badge{NCDNavy}{Sofie 1}\ Hej Peter! Jeg har desværre fået lidt ondt i mit højre knæ efter vores lange løbetur i søndags. Tror du, jeg kan være med til intervalløb i aften?\par\vspace{4pt}
\textbf{0. Peter:}\ \framebox[0.85cm]{\textbf{Z}} \quad \textcolor{NCDGray}{\small (Eksempel)} &
\badge{NCDGray}{Z}\ \textcolor{NCDGray}{\textit{Hej Sofie! Det er jeg ked af at høre. Hvis du har smerter, skal du helt sikkert holde en pause fra hårdt løb i aften.}} \\
\midrule
\badge{NCDNavy}{Sofie 2}\ Okay, det har du nok ret i, selvom det er irriterende at gå glip af træningen. Hvor længe skal jeg vente, før jeg må begynde at løbe igen?\par\vspace{4pt}
\textbf{1. Peter:}\ \framebox[0.85cm]{\rule{0pt}{0.26cm}} &
\badge{NCDNavy}{A}\ Ja, tag bare med til intervalløb i aften alligevel, så går smerterne nok væk af sig selv, når du bliver varm.\par\vspace{6pt}
\badge{NCDNavy}{B}\ Du skal helst sidde helt stille på sofaen i syv dage og undgå både vand og strækøvelser. \\
\midrule
\badge{NCDNavy}{Sofie 3}\ En hel uge uden løb lyder lidt langt, men jeg vil jo heller ikke have en alvorlig skade. Må jeg slet ikke bevæge mig i den uge, eller kan jeg træne noget andet i stedet?\par\vspace{4pt}
\textbf{2. Peter:}\ \framebox[0.85cm]{\rule{0pt}{0.26cm}} &
\badge{NCDNavy}{C}\ Start meget langsomt på fladt underlag, og husk at tjekke, om dine løbesko er for gamle og skal skiftes ud.\par\vspace{6pt}
\badge{NCDNavy}{D}\ Jeg vil anbefale dig at tage en uges pause fra løb, indtil knæet slet ikke gør ondt længere. \\
\midrule
\badge{NCDNavy}{Sofie 4}\ Svømning og udstrækning lyder faktisk som en god plan, og der ligger heldigvis en svømmehal tæt på mit arbejde. Er der noget særligt, jeg skal huske, før jeg prøver at løbe en kort tur igen?\par\vspace{4pt}
\textbf{3. Peter:}\ \framebox[0.85cm]{\rule{0pt}{0.26cm}} &
\badge{NCDNavy}{E}\ Jeg foreslår, at du løber en hurtig testtur på ti kilometer i morgen tidlig for at mærke, om knæet holder.\par\vspace{6pt}
\badge{NCDNavy}{F}\ Jo, du må meget gerne holde kroppen i gang. Prøv at svømme et par gange og lave rolige strækøvelser derhjemme, fordi det ikke belaster knæet. \\
\midrule
\badge{NCDNavy}{Sofie 5}\ Tak for de gode råd, Peter! Så køber jeg et par nye sko på fredag og starter stille ud i næste uge. Vi ses til træning snart! &
\textcolor{NCDGray}{\small （第 5 条为 Sofie 的收尾回复，无需补全）} \\
\end{longtable}
\normalsize

% ---------------------------------------------------------------------
\subsubsection*{Opgave 3 (NY · 5 point)：Amira starter i en madklub}

\textbf{Instruktion:}
\begin{itemize}
  \item Læs teksten \textit{Amira starter i en madklub}.
  \item Der er fem afsnit.
  \item I hvert afsnit mangler der en sætning.
  \item Under hvert afsnit står der fire sætninger.
  \item Skriv bogstavet på den sætning, der passer i afsnittet.\\
  Der er tre sætninger, du ikke skal bruge.
\end{itemize}

\vspace{0.25em}
\textbf{\color{NCDNavy}Amira starter i en madklub}

\vspace{0.2em}
\textbf{1.} Amira er for nylig flyttet til Odense for at begynde på et nyt arbejde som bioanalytiker. I starten kendte hun næsten ingen i byen, og efter arbejde sad hun ofte alene i sin lejlighed med en hurtig ret fra supermarkedet. \rule{2.4cm}{0.4pt} En dag så hun et opslag på det lokale bibliotek om en madklub, hvor beboere i kvarteret mødtes hver anden uge for at lave mad sammen. \quad \textbf{Svar:} \rule{1.0cm}{0.4pt}
\begin{enumerate}[label=\textbf{\Alph*.}, leftmargin=2.2em, itemsep=0.12em]
  \item Det var hun rigtig glad for, fordi hun foretrak at spise i ro og mag derhjemme.
  \item Hun savnede både at lave rigtig mad fra bunden og at lære nye mennesker at kende.
  \item Derfor inviterede hun sine mange gamle naboer i Odense på middag hver fredag.
  \item Men hun havde faktisk slet ingen tid til at gå på biblioteket eller lave mad.
\end{enumerate}

\vspace{0.25em}
\textbf{2.} Klubben bestod af otte mennesker i forskellige aldre, som alle boede i det samme kvarter. Første gang Amira kom til mødet, var hun en smule nervøs, fordi hun ikke vidste, om hun ville passe ind i gruppen. \rule{2.4cm}{0.4pt} Allerede efter ti minutter stod hun ved komfuret og grinede sammen med en ældre herre, der hedder Jens, mens de skar grøntsager til en stor suppe. \quad \textbf{Svar:} \rule{1.0cm}{0.4pt}
\begin{enumerate}[label=\textbf{\Alph*.}, leftmargin=2.2em, itemsep=0.12em]
  \item Heldigvis blev hun mødt med stor varme og fik straks en opgave i køkkenet.
  \item Desværre havde de andre medlemmer allerede spist, da hun nåede frem.
  \item Derfor besluttede hun at tage hjem igen med det samme uden at hilse på nogen.
  \item Men ingen af deltagerne havde lyst til at snakke med nye medlemmer.
\end{enumerate}

\vspace{0.25em}
\textbf{3.} I madklubben skiftes medlemmerne til at bestemme menuen og købe ind til aftenen. Da det blev Amiras tur i november, ville hun gerne vise de andre en traditionel ret fra sit hjemland. \rule{2.4cm}{0.4pt} Amira var lidt bange for, at krydderierne ville være for stærke for de danske deltagere, men gryderetten blev hurtigt spist op, og flere bad straks om opskriften. \quad \textbf{Svar:} \rule{1.0cm}{0.4pt}
\begin{enumerate}[label=\textbf{\Alph*.}, leftmargin=2.2em, itemsep=0.12em]
  \item Hun valgte derfor at bestille pizza udefra i stedet for at stå i køkkenet.
  \item De andre sagde dog nej, fordi de kun ønskede at spise klassiske danske retter.
  \item Hun foreslog dog først at lave mad fra et helt andet land, som hun ikke kendte.
  \item Derfor fandt hun en gammel familieopskrift frem og brugte hele eftermiddagen på at forberede krydderierne.
\end{enumerate}

\vspace{0.25em}
\textbf{4.} Ud over selve maden har aftenerne i klubben også haft en anden vigtig betydning for Amira. Når de otte deltagere sidder omkring bordet og taler om arbejde, nyheder og hverdagsliv, får hun trænet sit danske sprog i en afslappet stemning. \rule{2.4cm}{0.4pt} På den måde føler hun sig nu meget mere sikker, når hun skal diskutere eller småsnakke med kollegerne i laboratoriet. \quad \textbf{Svar:} \rule{1.0cm}{0.4pt}
\begin{enumerate}[label=\textbf{\Alph*.}, leftmargin=2.2em, itemsep=0.12em]
  \item Men de øvrige medlemmer vil desværre helst tale engelsk med hende hele aftenen.
  \item Desuden lærer hun mange nye hverdagsudtryk, som ikke står i hendes lærebøger.
  \item Derfor har hun valgt at droppe sit danskkursus og kun læse bøger derhjemme.
  \item Hun synes nemlig, at det er alt for kedeligt at snakke om hverdagen over maden.
\end{enumerate}

\vspace{0.25em}
\textbf{5.} Nu har Amira været medlem af madklubben i næsten et år, og tirsdag aften er blevet hendes yndlingsdag i kalenderen. Selvom hun stadig arbejder meget og nogle gange kommer træt hjem, melder hun sjældent afbud. \rule{2.4cm}{0.4pt} Faktisk er hun allerede i gang med at planlægge næste måneds menu sammen med Jens. \quad \textbf{Svar:} \rule{1.0cm}{0.4pt}
\begin{enumerate}[label=\textbf{\Alph*.}, leftmargin=2.2em, itemsep=0.12em]
  \item Hun regner dog med at melde sig ud næste måned for at få mere tid til sig selv.
  \item Hvis hun har været træt på arbejde, bliver hun næsten altid hjemme på sofaen.
  \item Hun oplever nemlig, at de fælles middage giver hende fornyet energi til resten af ugen.
  \item Jens og de andre har nemlig besluttet at lukke madklubben efter jul.
\end{enumerate}

% ---------------------------------------------------------------------
\subsubsection*{Opgave 4 (NY · 7 point)：Tre personer fortæller om deres forhold til mobil og skærmtid}

\textbf{Instruktion:}
\begin{itemize}
  \item Læs de tre tekster, og læs spørgsmålene (1--7).
  \item Find den person, der passer til hvert af de syv spørgsmål.
  \item Sæt kryds (\textbf{X}) ud for den person (A, B eller C), der passer til spørgsmålet. Se eksemplet (0).
\end{itemize}

\vspace{0.2em}
\small
\renewcommand{\arraystretch}{1.3}
\begin{longtable}{@{} >{\raggedright\arraybackslash}p{11.0cm} @{\hspace{0.3cm}} >{\centering\arraybackslash}p{1.65cm} @{\hspace{0.15cm}} >{\centering\arraybackslash}p{1.65cm} @{\hspace{0.15cm}} >{\centering\arraybackslash}p{1.65cm} @{}}
\toprule
\textbf{Spørgsmål} & \textbf{A. Ida} & \textbf{B. Tobias} & \textbf{C. Lin} \\
\midrule
\endfirsthead
\toprule
\textbf{Spørgsmål（续表）} & \textbf{A. Ida} & \textbf{B. Tobias} & \textbf{C. Lin} \\
\midrule
\endhead
\bottomrule
\endfoot
0. Hvem ser næsten ikke på en skærm i løbet af sin arbejdsdag? \textcolor{NCDGray}{\small (Eksempel)} & A & \textbf{X} & C \\
\midrule
1. Hvem lader ikke sin telefon ligge i soveværelset om natten længere? & A & B & C \\
\midrule
2. Hvem bruger mobilen til at øve sig i et fremmedsprog? & A & B & C \\
\midrule
3. Hvem opdagede sit store skærmforbrug, fordi et familiemedlem sagde noget om det? & A & B & C \\
\midrule
4. Hvem har en aftale med sine venner om at lægge telefonerne væk, når de ses? & A & B & C \\
\midrule
5. Hvem synes, at digitale enheder gør det lettere at bevare kontakten til folk langt væk? & A & B & C \\
\midrule
6. Hvem sætter tid af på én bestemt ugedag til helt at slukke for mobilen? & A & B & C \\
\midrule
7. Hvem sover bedre nu, efter at personen ændrede sine aftenregler? & A & B & C \\
\end{longtable}
\normalsize

\vspace{0.3em}
\begin{center}
\textbf{\large\color{NCDNavy} Tre personer fortæller om deres forhold til mobil og skærmtid}
\end{center}

\begin{emailbox}{A. Ida}
Jeg arbejder som projektkoordinator i et mellemstort firma i Aarhus, hvor jeg får beskeder på både computer og telefon hele dagen lang. Tidligere tog jeg ofte arbejdet med i seng, fordi jeg lige ville tjekke de sidste mails inden sengetid. Det ødelagde min nattesøvn fuldstændigt, for hovedet fortsatte med at tænke på opgaver hele natten. For et halvt år siden lavede jeg derfor faste aftenregler: Efter klokken nitten svarer jeg aldrig på jobmails, og min mobil bliver liggende ude i gangen hele natten, mens et gammelt vækkeur vækker mig om morgenen. Jeg falder meget hurtigere i søvn nu og vågner faktisk frisk og udhvilet. Derudover har mine veninder og jeg indført en fast aftale, når vi mødes på café efter arbejde: Alle telefoner bliver lagt midt på bordet på lydløs, indtil vi går hjem, så vi virkelig er til stede sammen.
\end{emailbox}

\begin{emailbox}{B. Tobias}
Som tømrer i Vejle arbejder jeg udendørs eller på værksted det meste af dagen, så jeg kigger næsten aldrig på en skærm i arbejdstiden. Til gengæld opdagede jeg sidste efterår, at jeg brugte over to timer hver eneste aften hjemme i sofaen på at se sportsvideoer og læse nyheder på nettet. Det var faktisk min tiårige søn, der fik mig til at tænke over det, da han en dag spurgte, hvorfor jeg altid kiggede ned i mobilen, når han ville snakke med mig. Den bemærkning ramte mig ret hårdt. Nu lægger jeg altid telefonen på en høj hylde i køkkenet fra klokken fem til otte om aftenen, så vi kan spise aftensmad og spille brætspil sammen i fred og ro. Desuden har jeg indført en fast mobilfri søndag, hvor min telefon er helt slukket fra morgen til aften, mens hele familien cykler eller fisker.
\end{emailbox}

\begin{emailbox}{C. Lin}
Jeg er medicinstuderende på universitetet i Odense og bruger digitale enheder rigtig mange timer i døgnet. Alle mine lærebøger og forelæsningsnoter ligger nemlig på min tablet, så jeg kan slet ikke undvære en skærm i hverdagen. Samtidig bruger jeg mobilen aktivt til at øve dansk, når jeg tager den lange bustur til campus hver morgen: Jeg lytter til danske nyhedspodcasts og træner nye ord i en sprogapp. Mange siger, at telefoner kun skaber afstand mellem mennesker, men for mig er det helt omvendt. Min nærmeste familie bor i Singapore, og fordi jeg ringer til mine forældre på video hver søndag og sender små billeder fra mit liv i Danmark, føler jeg mig tæt på dem på trods af den store afstand. Jeg tjekker stadig mine beskeder sent om aftenen i soveværelset på grund af tidsforskellen, uden at det skader mit studie næste morgen.
\end{emailbox}

% =====================================================================
% DEL 2: SKRIVNING
% =====================================================================
\ncdsubsec{Del 2: Skrivning（书面表达｜Opgave 1 Klage \& Opgave 2 E-mail）}

\begin{tipbox}{DU3 Modul 3 书面双任务考场提示}
\small
本部分包含两道写作任务：\textbf{Opgave 1}（正式/半正式投诉或推荐信，必须完整覆盖全部 4 个要点；建议写到 105--135 词并用 \textit{I / jer / jeres} 指代机构）；\textbf{Opgave 2}（回复朋友邮件，官方硬性规定 \textbf{Minimum 90 ord · Tid: 30 minutter}，需包含得体开头、结尾并回答全部提问）。
\end{tipbox}

\subsubsection*{Opgave 1（正式投诉信 · Klage til Hotel Kystblik）}

\textbf{Situation:}\\
Du boede i weekenden på \textit{Hotel Kystblik}, men opholdet var skuffende. Der var kraftig støj fra gaden hele natten, vandet i bruseren var koldt, rengøringen på værelset var dårlig, og receptionisten var uhøflig, da du sagde til.

\vspace{0.3em}
\textbf{Opgave:}\\
Skriv en klage til \textit{Hotel Kystblik}. Du skal fortælle:
\begin{itemize}
  \item hvornår du boede der
  \item hvilke problemer der opstod
  \item hvordan personalet reagerede
  \item hvad du forventer som kompensation
\end{itemize}

\begin{emailbox}{Din klage｜Til: Hotel Kystblik \qquad Emne: Klage over ophold}
\vspace{0.3em}
\rule{\linewidth}{0.4pt}\\[0.55em]
\rule{\linewidth}{0.4pt}\\[0.55em]
\rule{\linewidth}{0.4pt}\\[0.55em]
\rule{\linewidth}{0.4pt}\\[0.55em]
\rule{\linewidth}{0.4pt}\\[0.55em]
\rule{\linewidth}{0.4pt}\\[0.55em]
\rule{\linewidth}{0.4pt}
\end{emailbox}

\clearpage
\subsubsection*{Opgave 2（私人回信 · Minimum 90 ord · Tid: 30 minutter）}

\textbf{Situation:}\\
Du har fået e-mailen her nedenfor.

\vspace{0.3em}
\textbf{Opgave:}\\
Læs e-mailen og skriv et svar til Astrid. Du skal svare på alle Astrids spørgsmål. Du skal indlede og afslutte e-mailen på en passende måde.\\
\textbf{Længde:} Minimum 90 ord. \qquad \textbf{Tid:} 30 minutter.

\begin{emailbox}{Modtaget e-mail fra Astrid}
Hej!

Jeg hørte fra fælles venner, at du er startet på et spændende aftenkursus eller en ny fritidsinteresse efter arbejde! Det lyder virkelig sejt, men også lidt travlt. Hvilket kursus eller hvilken aktivitet går du til, og hvorfor valgte du det? Hvordan når du at balancere dit arbejde, din danskundervisning og din fritid i hverdagen? Og hvad gør du for at slappe af i weekenden, når ugen har været hård? Jeg glæder mig til at høre fra dig!

Kærlig hilsen\\
Astrid
\end{emailbox}

\begin{emailbox}{Dit svar｜Til: Astrid \qquad Emne: Sv: Aftenkursus og hverdag}
\vspace{0.3em}
\rule{\linewidth}{0.4pt}\\[0.55em]
\rule{\linewidth}{0.4pt}\\[0.55em]
\rule{\linewidth}{0.4pt}\\[0.55em]
\rule{\linewidth}{0.4pt}\\[0.55em]
\rule{\linewidth}{0.4pt}\\[0.55em]
\rule{\linewidth}{0.4pt}\\[0.55em]
\rule{\linewidth}{0.4pt}\\[0.55em]
\rule{\linewidth}{0.4pt}\\[0.55em]
\rule{\linewidth}{0.4pt}
\end{emailbox}

% =====================================================================
% DEL 3: MUNDTLIG KOMMUNIKATION
% =====================================================================
\ncdsubsec{Del 3: Mundtlig kommunikation（口语交际全真题卡｜Opgave 1 \& 2）}

\subsubsection*{Mundtlig Opgave 1 (Emneopgave med Mindmap｜Del 1: 1--2 min. + Del 2: 3--4 min.)}

\begin{tipbox}{考场规则提醒（DU3 Modul 3 Opgave 1 · 难度与关键词红线）}
\small
考生需自带\textbf{两张}提前填写好的思维导图（至少涉及工作/学习/医疗健康/环保/数字化/家务分工等达到 B1 复杂度的主题）。导图框内\textbf{只能写关键词或短语（stikord），严禁写完整句子}，以防考场照本宣科！抽签选定 1 张后，先进行 \textbf{1--2 分钟}连贯陈述（Del 1），随后回答考官 \textbf{3--4 分钟}追问（Del 2）并简要说明原因。\\
\textbf{官方禁用红线}：\textit{Min bolig/lejlighed}、\textit{Min familie} 与 \textit{Min hverdag} 复杂度停留在 A2，\textbf{不可作为 Modul 3 Opgave 1 话题}！
\end{tipbox}

\begin{grammatikbox}{Mindmap 2A（已填关键词卡）：Mine sunde og usunde vaner}
\small
\renewcommand{\arraystretch}{1.35}
\begin{tabularx}{\linewidth}{@{} >{\leavevmode\bfseries\color{NCDNavy}\raggedright\arraybackslash}p{3.1cm} @{\hspace{0.2cm}} >{\raggedright\arraybackslash}X @{\hspace{0.35cm}} >{\leavevmode\bfseries\color{NCDNavy}\raggedright\arraybackslash}p{3.1cm} @{\hspace{0.2cm}} >{\raggedright\arraybackslash}X @{}}
\toprule
\multicolumn{4}{c}{\large\bfseries\color{NCDRed} Emne: Mine sunde og usunde vaner \quad\small\color{NCDGray}（仅含 stikord 关键词，不含整句）} \\
\midrule
1. sund mad & havregryn · frisk frugt og grøntsager · fuldkornsbrød · fisk 2 gange ugentligt · meget vand &
2. usund mad & chokolade om aftenen · fredagsslik · fastfood ved travlhed · op til 4 kopper kaffe \\
\midrule
3. motion & cykle 8 km til arbejde · svømmehal onsdag aften · lang gåtur i skoven om søndagen &
4. rygning / alkohol & aldrig røget · ingen cigaretter · 1 glas vin i weekenden · alkoholfrit øl til hverdag \\
\midrule
5. søvn & ca. 7 timers søvn · fast sengetid kl. 23 · ingen mobil efter kl. 22 · træt mandag morgen &
6. mening om din livsstil (kan godt lide / ikke lide) & tilfreds med motion og mad · god balance · vil gerne skære ned på sukker · ikke for streng \\
\bottomrule
\end{tabularx}
\end{grammatikbox}

\textbf{Del 2 考官追问模拟（Mine sunde og usunde vaner · 4 题）：}
\begin{enumerate}
  \item \textit{Du siger, at du spiser lidt for meget chokolade om aftenen -- hvornår har du mest lyst til usunde ting, og hvorfor?}
  \item \textit{Hvorfor kan du godt lide at svømme og cykle til arbejde, og hvordan holder du motivationen om vinteren?}
  \item \textit{Du nævner, at du lægger telefonen væk klokken ti -- har det gjort en forskel for din søvn?}
  \item \textit{Hvad synes du generelt er det vigtigste for at leve et sundt liv i en travl hverdag i Danmark?}
\end{enumerate}

\begin{grammatikbox}{Mindmap 2B（已填关键词卡）：Mine digitale vaner}
\small
\renewcommand{\arraystretch}{1.35}
\begin{tabularx}{\linewidth}{@{} >{\leavevmode\bfseries\color{NCDNavy}\raggedright\arraybackslash}p{3.1cm} @{\hspace{0.2cm}} >{\raggedright\arraybackslash}X @{\hspace{0.35cm}} >{\leavevmode\bfseries\color{NCDNavy}\raggedright\arraybackslash}p{3.1cm} @{\hspace{0.2cm}} >{\raggedright\arraybackslash}X @{}}
\toprule
\multicolumn{4}{c}{\large\bfseries\color{NCDRed} Emne: Mine digitale vaner \quad\small\color{NCDGray}（仅含 stikord 关键词，不含整句）} \\
\midrule
1. skærmtid / telefon & ca. 2,5 timer dagligt · tidsgrænse på apps · telefon i tasken under aftensmad &
2. sociale medier & LinkedIn og lidt Instagram · slået notifikationer fra · sjældent opslag selv \\
\midrule
3. nyheder / podcasts & DR Nyheder om morgenen · danske podcasts på cyklen · øve dansk udtale &
4. nethandel / MitID & MitID til netbank, e-Boks og borger.dk · undgå impulskøb · lokale butikker \\
\midrule
5. arbejde / studie online & onlinemøder i Teams · digitale lærebøger og ordbøger · huske korte skærmpauser &
6. mening om digitale vaner (kan godt lide / ikke lide) & praktisk og tidsbesparende · nem kontakt · bekymret for uro · ønsker flere skærmfrie aftener \\
\bottomrule
\end{tabularx}
\end{grammatikbox}

\textbf{Del 2 考官追问模拟（Mine digitale vaner · 4 题）：}
\begin{enumerate}
  \item \textit{Du fortæller, at du har slået notifikationer fra -- hvordan har det påvirket din hverdag på arbejde?}
  \item \textit{Hvad synes du om, at næsten alting i Danmark foregår digitalt med MitID og e-Boks?}
  \item \textit{Vil du fortælle lidt mere om, hvordan du bruger danske podcasts eller apps til at lære dansk?}
  \item \textit{Tror du, vi bruger telefonen for meget i dag, eller er der mest fordele ved digitale løsninger?}
\end{enumerate}

% ---------------------------------------------------------------------
\subsubsection*{Mundtlig Opgave 2 (Paropgave Del 1 ca. 4 min. + Individuel Del 2 ca. 3 min. pr. kursist)}

\begin{center}
\large\bfseries\color{NCDTeal} OPGAVEARK \qquad|\qquad EMNE: Fritid og motion
\end{center}

\small
\renewcommand{\arraystretch}{1.35}
\begin{tabularx}{\linewidth}{@{} >{\raggedright\arraybackslash}X @{\hspace{0.4cm}} >{\raggedright\arraybackslash}X @{}}
\toprule
\badge{NCDNavy}{Billede 1}\ \textbf{\color{NCDNavy}[图片 1: 健身房器械训练]}\par\vspace{3pt}
\textit{Træning i et moderne fitnesscenter med vægte, løbebånd og maskiner indendørs.}
&
\badge{NCDNavy}{Billede 2}\ \textbf{\color{NCDNavy}[图片 2: 森林户外跑步]}\par\vspace{3pt}
\textit{Løbetur alene eller med en ven på en grøn skovsti i frisk luft.}
\\
\midrule
\badge{NCDNavy}{Billede 3}\ \textbf{\color{NCDNavy}[图片 3: 羽毛球/排球等团体球类]}\par\vspace{3pt}
\textit{Holdsport i en lokal sportshal, hvor man spiller badminton eller volleyball sammen.}
&
\badge{NCDNavy}{Billede 4}\ \textbf{\color{NCDNavy}[图片 4: 居家瑜伽与拉伸]}\par\vspace{3pt}
\textit{Rolig yoga og udstrækning på en træningsmåtte i stuen derhjemme.}
\\
\bottomrule
\end{tabularx}
\normalsize

\vspace{0.35em}
\begin{tipbox}{TESTERARK｜考官官方引导语与 Del 1 / Del 2 考核流程}
\small
\textbf{Del 1 (Paropgave -- ca. 4 min.)}\\
Opgavearket lægges foran kursisterne, der præsenteres for emnet på følgende måde:\\[3pt]
\textit{Emnet er 'Fritid og motion'. Her er nogle billeder, der viser forskellige måder at dyrke motion på i sin fritid. I skal tale med hinanden om, hvad I synes om de forskellige måder at dyrke motion på, og hvorfor. I skal stille spørgsmål til hinanden og svare på hinandens spørgsmål.}\\[4pt]
Hvis kursisterne ikke udtrykker deres præferencer i forhold til alle fire billeder, må tester sørge for at dette sker ved at spørge:\\[2pt]
\textit{Kan I lide at træne i et fitnesscenter / løbe ude i naturen / dyrke holdsport i en hal / lave yoga og udstrækning derhjemme? Hvorfor? / Hvorfor ikke?}\\[6pt]
\textbf{Del 2 (Individuel opgave -- ca. 3 min. pr. kursist)}\\
Tester stiller hver kursist nedenstående spørgsmål:
\begin{enumerate}[leftmargin=1.6em, itemsep=0.15em]
  \item \textit{Hvordan kan du bedst lide at slappe af eller dyrke motion i din fritid? Hvorfor?}
  \item \textit{Hvornår og hvor ofte har du tid til fritidsinteresser i en almindelig uge?}
  \item \textit{Kan du fortælle lidt om en fritidsaktivitet eller sport, du har prøvet eller gerne vil prøve? (hvor, med hvem, hvorfor, nemt/svært)}
\end{enumerate}
Hvis kursisten har besvaret et spørgsmål under samtalen i del 1, stilles uddybende spørgsmål, fx: \textit{Du siger, at \dots{} vil du fortælle lidt mere om \dots{}.} Der stilles evt. opfølgende og uddybende spørgsmål inden for tidsrammen.
\end{tipbox}

% =====================================================================
% DEL 4: RETTEARK, DETALJERET ANALYSE & HØJSCORE-MODELLER
% =====================================================================
\ncdsubsec{Del 4: Retteark, Detaljeret Analyse \& Højscore-Modeller}

\subsubsection*{4.1 官方 Retteark 评分卡（Læsning 24 point）}

\small
\renewcommand{\arraystretch}{1.28}
\begin{longtable}{@{} >{\leavevmode\bfseries\color{NCDNavy}\raggedright\arraybackslash}p{2.9cm} @{\hspace{0.25cm}} >{\raggedright\arraybackslash}p{9.6cm} @{\hspace{0.25cm}} >{\centering\arraybackslash}p{1.75cm} @{\hspace{0.25cm}} >{\centering\arraybackslash}p{1.65cm} @{}}
\toprule
\textbf{模块 / 题号} & \textbf{Rettenøgle（标准答案）} & \textbf{分值} & \textbf{自测得分} \\
\midrule
\endfirsthead
\toprule
\textbf{模块 / 题号（续表）} & \textbf{Rettenøgle（标准答案）} & \textbf{分值} & \textbf{自测得分} \\
\midrule
\endhead
\bottomrule
\endfoot
\textbf{Læsning Opgave 1}\newline（6 题） &
1. \textbf{før} \quad 2. \textbf{fordi} \quad 3. \textbf{dog} \quad 4. \textbf{til} \quad 5. \textbf{at} \quad 6. \textbf{givet}\newline
\textcolor{NCDGray}{\small （多余干扰词：blev, eller, fra, med, men, som｜例词 0 为 det）} & 6 point & \rule{0.75cm}{0.4pt} / 6 \\
\midrule
\textbf{Læsning Opgave 2A}\newline（3 题 · 聊天补全） &
1. \textbf{B} \quad 2. \textbf{E} \quad 3. \textbf{A}\newline
\textcolor{NCDGray}{\small （多余干扰项：C, D, F｜例句 0 为 Z）} & 3 point & \rule{0.75cm}{0.4pt} / 3 \\
\midrule
\textbf{Læsning Opgave 2B}\newline（3 题 · 聊天补全） &
1. \textbf{D} \quad 2. \textbf{F} \quad 3. \textbf{C}\newline
\textcolor{NCDGray}{\small （多余干扰项：A, B, E｜例句 0 为 Z）} & 3 point & \rule{0.75cm}{0.4pt} / 3 \\
\midrule
\textbf{Læsning Opgave 3}\newline（5 题 · 段内插句） &
1. \textbf{B} \quad 2. \textbf{A} \quad 3. \textbf{D} \quad 4. \textbf{B} \quad 5. \textbf{C} & 5 point & \rule{0.75cm}{0.4pt} / 5 \\
\midrule
\textbf{Læsning Opgave 4}\newline（7 题 · 三人匹配） &
1. \textbf{A} \quad 2. \textbf{C} \quad 3. \textbf{B} \quad 4. \textbf{A} \quad 5. \textbf{C} \quad 6. \textbf{B} \quad 7. \textbf{A}\newline
\textcolor{NCDGray}{\small （例题 0 为 B. Tobias）} & 7 point & \rule{0.75cm}{0.4pt} / 7 \\
\midrule
\textbf{Læsning 总计} & \textbf{Bestået：15--24 point} \quad|\quad \textbf{Ikke bestået：0--14 point} & \textbf{24 point} & \rule{0.75cm}{0.4pt} / 24 \\
\end{longtable}
\normalsize

\begin{tipbox}{Skrivning \& Mundtlig 自检过关准则（Helhedsvurdering）}
\small
\begin{itemize}
  \item \textbf{Skrivning（整体评估 Bestået）}：\ding{172} Opgave 1 投诉信必须不漏任何一条指定要点，且全程用 \textbf{\textit{I / jer / jeres}} 称呼酒店（切忌中途跳成 \textit{du / dig}）；\ding{173} Opgave 2 私人邮件严格满 \textbf{90 词以上}（建议写到 105--135 词），回答 Astrid 全部 4 问，并含得体称呼与落款；\ding{174} 主动运用 NCD 3 核心语法（\textit{-s/blive} 被动、\textit{til trods for at / på baggrund af}、从句中否定副词前置等）。
  \item \textbf{Mundtlig Opgave 1 \& 2（整体评估 Bestået）}：Opgave 1 陈述紧扣 6 个 \textit{stikord} 并简要说明原因；Opgave 2 双人对话必须对\textbf{全部 4 张运动图}明确表达喜欢/不喜欢并给出理由，同时积极向搭档发问。
\end{itemize}
\end{tipbox}

% ---------------------------------------------------------------------
\subsubsection*{4.2 阅读 Opgave 1--4 逐题解析、全文精译与核心词表}

\paragraph*{一、Opgave 1（一周无社交媒体实验）词汇语法定位与干扰项排除}
\begin{itemize}
  \item \textbf{(0) det [Eksempel]}：回指前分句“每天花近三小时刷社交媒体”这件事，作 \textit{begyndte} 的形式/回指主语。
  \item \textbf{1. før}：时间从句连词“在……之前”，与 \textit{overhovedet stod ud af sengen} 搭配（“甚至还没下床就立刻查看消息”）。\textit{men / eller / som} 无法引导该时间状语从句，\textit{fra} 与 \textit{overhovedet} 语义不合。
  \item \textbf{2. fordi}：引导原因状语从句；注意从句内部状语 \textit{hele tiden} 位于变位动词 \textit{lyste} \textbf{之前}（典型从句语序），由此直接排除并列连词 \textit{men / eller}（并列句不改变常规主句语序）及介词/关系代词 \textit{fra / med / som}。
  \item \textbf{3. dog}：句中让步对比副词“然而/不过”，承接上文“前两天出奇难熬”与后文“三天后一切\textbf{却}平静得多”。该位置位于主句动词 \textit{blev} 与表语 \textit{meget roligere} 之间，其余候选词均不符合词性或句法位置。
  \item \textbf{4. til}：高频固定搭配 \textit{tid til at + 不定式}（“有时间去做某事”），不可替换为 \textit{med / fra}。
  \item \textbf{5. at}：引导 \textit{mærkede også tydeligt} 的宾语从句（“明显感觉到……”），只有连接词 \textit{at} 合法。
  \item \textbf{6. givet}：跟在助动词 \textit{har} 后构成现在完成时，且后面紧接双宾语 \textit{ham mere energi}（“给了他更多精力”）；其余干扰词 \textit{blev} 不能接在 \textit{har} 后，也无法带双宾语。
\end{itemize}

\begin{transbox}
\small
\textbf{\color{NCDTeal}[Opgave 1 全文中文精译]}\\
米克尔今年二十九岁，在罗斯基勒担任平面设计师。多年来，他每天几乎花三个小时在社交媒体上，\textbf{这件事}慢慢开始让他感到烦躁。每天早上闹钟一响，他甚至\textbf{在}还没下床\textbf{之前}就立刻查看消息。在工作时他也常常难以集中注意力，\textbf{因为}屏幕总是不断亮起新通知和短视频。于是上个月他决定做一个实验：把手机上所有社交软件删除七天，感受一下变化。最初两天出奇地难熬。每当米克尔坐在火车上或在食堂等咖啡时，手几乎会条件反射般地伸进口袋掏手机。\textbf{不过}三天之后，一切都变得平静多了。少了源源不断的图片信息流，他晚上睡得更香，还突然有了时间\textbf{去}读一本在床头柜躺了半年的小说。他也明显感觉到，\textbf{（就是）}自己在休息时间与同事聊得更多了，而不是只低头盯着屏幕。等这周终于结束时，米克尔只重新下载了两个应用，并且关掉了所有推送通知。现在他每天最多只花二十分钟上网，而这\textbf{给}他在日常生活中带来了更多精力。
\end{transbox}

\paragraph*{二、Opgave 2A \& 2B（聊天对话补全）双向衔接解析与精译}
\begin{itemize}
  \item \textbf{Opgave 2A（Lukas 与 Emma 合买周年礼物）}：
    \begin{itemize}
      \item \textbf{1. $\rightarrow$ B}：Lukas 2 询问每人分摊多少（\textit{Hvor meget... hver person}），B 提出“每人 100 克朗”且顺带引出“还需要买配咖啡的点心（\textit{noget lækkert til kaffen}）”，正好触发 Lukas 3 的“8 人共 800 克朗”及“要不要我顺便去面包店订蛋糕”。干扰项 D 虽出现 \textit{100 kroner hver}，却否定收钱，与 Lukas 3 的 \textit{i alt 800 kroner} 直接矛盾。
      \item \textbf{2. $\rightarrow$ E}：Lukas 3 主动问能否去订蛋糕，E 欣然同意并请他预订“无坚果巧克力蛋糕（\textit{chokoladekage uden nødder}）”，紧接 Lukas 4 的“好的，那我明天上午打电话订无坚果巧克力蛋糕”。干扰项 C 说“根本不需要蛋糕”，与 Lukas 4 矛盾。
      \item \textbf{3. $\rightarrow$ A}：Lukas 4 问两人是否知道周六聚会，A 回答“他们毫不知情，哥哥 14:30 悄悄放我们进去”，完美引出 Lukas 5 的“太好了，那我当然保密，周六 14:30 见”。干扰项 F 说“已发邀请、15 点见”，与保密及时间双双冲突。
    \end{itemize}
  \item \textbf{Opgave 2B（Sofie 与跑步教练 Peter 沟通膝伤）}：
    \begin{itemize}
      \item \textbf{1. $\rightarrow$ D}：Sofie 2 问需要停跑多久，D 建议“停跑一周直到膝盖完全不痛”，正对 Sofie 3 的“一整周不跑听起来有点久（\textit{En hel uge uden løb}）”。干扰项 A 劝她今晚照常间歇跑，与例句 Z 矛盾。
      \item \textbf{2. $\rightarrow$ F}：Sofie 3 问这周能不能改练别的项目，F 建议“游泳和居家轻柔拉伸（\textit{svømme ... rolige strækøvelser}）”，正对 Sofie 4 开头的“游泳和拉伸听起来确实是好计划”。干扰项 B 说“七天一动不动并避开水和拉伸”，与 Sofie 4 冲突。
      \item \textbf{3. $\rightarrow$ C}：Sofie 4 问复跑前有什么注意事项，C 叮嘱“在平坦路面慢跑起步，并检查跑鞋是否太旧该换了”，正好引出 Sofie 5 的“那我周五就去买双新鞋，下周慢慢开始”。干扰项 E 建议“明早快跑 10 公里测试”，与休息一周和慢慢重启双双矛盾。
    \end{itemize}
\end{itemize}

\begin{transbox}
\small
\textbf{\color{NCDTeal}[Opgave 2A \& 2B 完整对话中文精译]}\\
\textbf{2A（Lukas \& Emma）}：L1：嗨艾玛！西格妮和卡斯珀下周六就是五周年纪念日了，你想好大家合送什么礼物了吗？ $\rightarrow$ \textbf{Z（例）}：嗨！还没完全定，不过我在想可以送一家好餐厅的礼品卡，你觉得呢？ $\rightarrow$ L2：好主意！他们俩都爱吃意餐。你觉得每人出多少钱合适？ $\rightarrow$ \textbf{1. B}：每人 100 克朗怎么样？谁都不会有压力，而且配咖啡我们也得准备点甜点。 $\rightarrow$ L3：很公平，那 8 个人总共就是 800 克朗。蛋糕也要我去面包店订吗？ $\rightarrow$ \textbf{2. E}：对，麻烦你了，因为你就住在面包店旁边。你能订一个不含坚果的巧克力蛋糕吗？ $\rightarrow$ L4：没问题，我明天上午就打电话订无坚果巧克力蛋糕。惊喜安排呢——他们知道我们周六下午过去吗？ $\rightarrow$ \textbf{3. A}：不知道，他们毫不知情！卡斯珀的弟弟会在 14:30 趁他们出门买东西时悄悄放我们进去。 $\rightarrow$ L5：太棒了，我肯定严格保密，这就给其他客人发集合时间。周六 14:30 见！\\[4pt]
\textbf{2B（Sofie \& Peter）}：S1：嗨彼得！周日长跑后我右膝盖有点疼，你觉得我今晚能参加间歇跑吗？ $\rightarrow$ \textbf{Z（例）}：嗨索菲！真遗憾听到这个消息。如果有疼痛，今晚绝对要暂停高强度跑步。 $\rightarrow$ S2：好吧，你说得对，虽然错过训练有点郁闷。我要等多久才能重新开始跑？ $\rightarrow$ \textbf{1. D}：我建议你暂停跑步一周，直到膝盖彻底不疼为止。 $\rightarrow$ S3：一周不跑听起来有点久，但我也不想落下重伤。这周我是一点都不能动，还是可以练点别的？ $\rightarrow$ \textbf{2. F}：可以保持身体活动。试着去游几次泳并在家做轻柔拉伸，因为这不会给膝盖施压。 $\rightarrow$ S4：游泳和拉伸听起来确实不错，而且我公司旁边刚好有游泳馆。试着短跑前有什么要特别记住的吗？ $\rightarrow$ \textbf{3. C}：一定在平坦路面上极慢地起步，而且记得检查跑鞋是不是太旧该换新了。 $\rightarrow$ S5：谢谢你的好建议，彼得！那我周五就去买双新鞋，下周轻轻恢复。训练场很快见！
\end{transbox}

\paragraph*{三、Opgave 3（Amira 加入美食俱乐部）段内插句精析与全文精译}
\begin{itemize}
  \item \textbf{第 1 段 $\rightarrow$ B}：前句写她刚搬来人生地不熟、晚饭总靠超市快餐将就；后句写她在图书馆看到“邻里做饭俱乐部”海报。选项 B“她既怀念亲手从头做热饭，也渴望认识新朋友”承前启后，直接解释了她的动机；A（更喜欢独自在家吃）、C（老邻居众多）、D（没时间去图书馆）均与上下文自相矛盾。
  \item \textbf{第 2 段 $\rightarrow$ A}：前句写她第一次赴约有点紧张（\textit{en smule nervøs}），担心融不进去；后句却画面一转“短短十分钟后她就和长者延斯在炉灶边有说有笑地切菜”。A“幸好她受到热情欢迎并立刻分到厨房任务”精准完成转折衔接；B（别人早吃完了）、C（立刻转身回家）、D（没人愿意和新成员说话）皆违背下文事实。
  \item \textbf{第 3 段 $\rightarrow$ D}：前句说轮到阿米拉定菜单，她想展示家乡传统菜；后句说她有点担心香料对丹麦人太辣，但炖菜很快被抢光。D“因此她翻出一份古老的家族食谱，花了一整个下午准备香料”将“传统菜”和后文的 \textit{krydderierne / gryderetten} 紧密扣合；A（点披萨）、B（大家只吃丹麦菜）、C（做别国菜）均断链。
  \item \textbf{第 4 段 $\rightarrow$ B}：前句指大家围坐聊天帮她练习了丹麦语；后句写“通过这种方式，她现在在实验室与同事闲聊讨论更有自信”。B“此外，她学到许多课本里找不到的日常地道表达”恰当延伸口语收获并顺接 \textit{På den måde}；其余选项（只说英语、退掉语言课、嫌聊天无聊）均错误。
  \item \textbf{第 5 段 $\rightarrow$ C}：前句说尽管工作累她也很少请假缺席（\textit{melder sjældent afbud}）；后句说她已经在和延斯策划下个月菜单。C“她感觉到共同晚餐能给剩下的一周注入新能量”以 \textit{nemlig} 解释原因并自然过渡；A（想退社）、B（累了就宅在沙发）、D（俱乐部圣诞后关闭）均与前后句冲突。
\end{itemize}

\begin{transbox}
\small
\textbf{\color{NCDTeal}[Opgave 3 全文中文精译]}\\
\textbf{1.} 阿米拉最近搬到欧登塞，开始担任生物分析员的新工作。起初她在城里几乎谁也不认识，下班后经常独自坐在公寓里对付一份超市快餐。\textbf{[B] 她既怀念亲手从头做一顿正餐，也渴望结识新朋友。}有一天，她在当地图书馆看到一张社区美食俱乐部的海报，邻里们每两周聚在一起做一次饭。\textbf{2.} 俱乐部由住在同一街区、年龄各异的八个人组成。阿米拉第一次去聚会时有点紧张，因为不知道自己能否融入这个集体。\textbf{[A] 幸好大家非常热情地欢迎了她，并立刻给她安排了厨房里的活儿。}短短十分钟后，她就已经站在炉灶旁，和一位名叫延斯的长者一边切汤菜一边开心大笑了。\textbf{3.} 在美食俱乐部里，成员们轮流决定菜单并采购食材。十一月轮到阿米拉做东时，她很想向大家展示一道家乡传统菜。\textbf{[D] 因此她找出一份古老的家族食谱，花了整个下午准备香料。}阿米拉原本有点担心香料对丹麦成员来说太辣，但那锅炖菜很快一扫而空，好几位成员立刻找她要食谱。\textbf{4.} 除了美食本身，俱乐部的夜晚对阿米拉还有另一重重要意义。当八位成员围坐在餐桌旁聊工作、新闻与日常生活时，她在轻松氛围中锻炼了丹麦语。\textbf{[B] 此外，她还学到了许多教材里没有的日常表达。}通过这种方式，她现在在实验室里与同事讨论或闲聊时感觉自信多了。\textbf{5.} 如今阿米拉加入俱乐部快一年了，周二晚上成了她日历上最喜欢的一天。尽管工作依然很忙、有时下班很累，她却很少请假缺席。\textbf{[C] 因为她感受到，一起共进晚餐能给她剩下的一周重新注入活力。}事实上，她已经在和延斯一起策划下个月的菜单了。
\end{transbox}

\paragraph*{四、Opgave 4（三人谈手机与屏幕时间）同义转述定位与精译}
\begin{itemize}
  \item \textbf{(0) B. Tobias [Eksempel]}：\textit{ser næsten ikke på en skærm i løbet af sin arbejdsdag} $\leftrightarrow$ \textit{kigger næsten aldrig på en skærm i arbejdstiden}（木匠工作几乎不看屏幕；Ida 全天用电脑手机回消息，Lin 白天用平板读书）。
  \item \textbf{1. A. Ida}：\textit{lader ikke sin telefon ligge i soveværelset om natten længere} $\leftrightarrow$ \textit{min mobil bliver liggende ude i gangen hele natten}（过去曾带手机上床，如今手机整晚留在走廊并换用老式闹钟；Lin 睡前仍在卧室刷消息）。
  \item \textbf{2. C. Lin}：\textit{øve sig i et fremmedsprog} $\leftrightarrow$ \textit{øve dansk ... lytter til danske nyhedspodcasts og træner nye ord i en sprogapp}。
  \item \textbf{3. B. Tobias}：\textit{fordi et familiemedlem sagde noget om det} $\leftrightarrow$ \textit{min tiårige søn, der fik mig til at tænke over det, da han en dag spurgte, hvorfor jeg altid kiggede ned i mobilen}（注意不要被 Ida 自己因失眠立规矩干扰）。
  \item \textbf{4. A. Ida}：\textit{aftale med sine venner om at lægge telefonerne væk} $\leftrightarrow$ \textit{mine veninder og jeg indført en fast aftale ... Alle telefoner bliver lagt midt på bordet på lydløs}。
  \item \textbf{5. C. Lin}：\textit{bevare kontakten til folk langt væk} $\leftrightarrow$ \textit{ringer til mine forældre på video hver søndag ... føler jeg mig tæt på dem på trods af den store afstand}。
  \item \textbf{6. B. Tobias}：\textit{sætter tid af på én bestemt ugedag til helt at slukke for mobilen} $\leftrightarrow$ \textit{indført en fast mobilfri søndag, hvor min telefon er helt slukket fra morgen til aften}（Ida 只是每晚不看工作邮件并放走廊，Lin 每周日反而要跟父母视频）。
  \item \textbf{7. A. Ida}：\textit{sover bedre nu, efter at personen ændrede sine aftenregler} $\leftrightarrow$ \textit{falder meget hurtigere i søvn nu og vågner faktisk frisk og udhvilet}（Tobias 仅谈及陪儿子玩桌游，Lin 从未调整晚间习惯）。
\end{itemize}

\begin{transbox}
\small
\textbf{\color{NCDTeal}[Opgave 4 三人自述中文精译]}\\
\textbf{A. 艾达}：我在奥胡斯一家中型公司做项目协调员，整天都要在电脑和手机上接收各种消息。以前我经常把工作带到床上，因为睡前总想再看一眼最后的邮件。这彻底毁了我的夜间睡眠，因为大脑整晚都在继续琢磨工作任务。因此半年前我制定了固定的晚间规定：晚上七点后绝不回复工作邮件，手机整晚放在门厅走廊里，早上靠一个老式闹钟叫醒我。现在我入睡快得多，醒来也真正感到神清气爽。此外，我和闺蜜们下班后约咖啡时也立下规矩：所有手机静音放在桌子正中央直到散场回家，这样大家才算真正专心相伴。\\
\textbf{B. 托比亚斯}：作为瓦埃勒的木匠，我白天大部分时间都在户外或工坊干活，所以工作时间几乎从不看屏幕。可另一方面，去年秋天我却发现自己每晚窝在沙发上刷体育视频、看网上新闻超过两个小时。其实是我十岁的儿子点醒了我——有一天他问我，为什么每次他想和我说话时，我总是在低头看手机。那句话深深戳中了我。现在每天傍晚五点到八点，我都会把手机放在厨房高处的架子上，好让我们安安静静地一起吃晚饭和玩桌游。此外我还推行了“无手机星期日”，周日从早到晚彻底关机，全家人一起骑车或钓鱼。\\
\textbf{C. 林}：我是欧登塞大学的医学生，一昼夜里使用电子设备的时间非常长。我的所有教材和讲义笔记都在平板电脑上，所以在日常生活中我根本离不开屏幕。与此同时，每天早上坐长途公交去校园时，我也会积极用手机练习丹麦语：我会听丹麦语新闻播客，并在语言软件里背新词。很多人说手机只会拉远人与人的距离，但对我来说恰恰相反。我的直系亲属住在新加坡，正因为我每周日都能跟父母视频通话，日常也能发些丹麦生活照，尽管相隔万里，我依然觉得与他们很亲近。由于时差，我睡前仍会在卧室看一眼消息，而这并不会影响我第二天早上的学业。
\end{transbox}

\paragraph*{五、核心高频词表（Gloser｜15 个必背屈折变化词条）}

\begin{glosertable}
名词 & en skærmtid & \textit{(skærmtiden, ingen flertal)} & 屏幕使用时间 \\
名词 & et eksperiment & \textit{(eksperimentet, eksperimenter, eksperimenterne)} & 实验，尝试 \\
动词 & at irritere & \textit{(irriterer, irriterede, har irriteret)} & 使烦躁，惹恼 \\
动词 & at stirre & \textit{(stirrer, stirrede, har stirret)} & 盯着看，凝视 \\
名词 & et jubilæum & \textit{(jubilæet, jubilæer, jubilæerne)} & 周年纪念日 \\
动词 & at anbefale & \textit{(anbefaler, anbefalede, har anbefalet)} & 推荐，建议 \\
动词 & at belaste & \textit{(belaster, belastede, har belastet)} & 给……施加负荷/压力 \\
短语 & at melde afbud & \textit{(melder afbud, meldte afbud, har meldt afbud)} & 取消预约，请假缺席 \\
动词 & at undvære & \textit{(undværer, undværede, har undværet)} & 缺少，离不开（\textit{ikke kunne undvære}） \\
副词 & til gengæld & （副词短语，句首触发 V2 倒装） & 相反，另一方面 \\
名词 & et ophold & \textit{(opholdet, ophold, opholdene)} & 逗留，住宿期间 \\
形容词 & uhøflig & \textit{(uhøflig, uhøfligt, uhøflige)} & 无礼的，态度恶劣的 \\
动词 & at refundere & \textit{(refunderer, refunderede, har refunderet)} & 退款，退还费用 \\
名词 & en balance & \textit{(balancen, balancer, balancerne)} & 平衡，均衡 \\
短语 & at koble af & \textit{(kobler af, koblede af, har koblet af)} & 彻底放松，摆脱忙碌 \\
\end{glosertable}

% =====================================================================
% 4.3 写作双档高分示范与点评
% =====================================================================
\subsubsection*{4.3 写作双档高分示范与点评（Skrivning Modeller）}

\begin{emailbox}{Opgave 1 范文一：B1 稳过版（115 词｜全覆 4 问 · 称谓一致 · 因果清晰）}
Kære Hotel Kystblik

Jeg skriver for at klage over mit ophold på jeres hotel fra den 14. til den 16. november, hvor jeg boede på værelse 204.

Desværre opstod der flere problemer. For det første var der kraftig støj fra gaden hele natten, så jeg næsten ikke sov. Derudover var vandet i bruseren helt koldt lørdag morgen, og værelset var ikke gjort ordentligt rent.

Da jeg gik ned i receptionen for at fortælle om dette, var receptionisten meget uhøflig. Han sagde bare kort, at han ikke kunne hjælpe, og gav ingen undskyldning.

Derfor forventer jeg, at I betaler halvdelen af prisen tilbage som kompensation. Jeg hører gerne fra jer hurtigst muligt.

Venlig hilsen\\
Katrine Holm

\vspace{0.3em}
\rule{\linewidth}{0.4pt}\\[0.2em]
\small\textbf{\color{NCDTeal}[考官点评 · 115 ord]}：稳健达标并覆盖全部 4 项要点（入住日期与房号、三大问题、前台反应、明确半价赔偿要求）。全文自始至终保持 \textit{jeres / I / jer} 第二人称复数指代酒店，结构连接词（\textit{For det første / Derudover / Da / Derfor}）运用地道，从句否定语序正确（\textit{så jeg næsten ikke sov}、\textit{at han ikke kunne hjælpe}）。
\end{emailbox}

\begin{emailbox}{Opgave 1 范文二：B1+ 高分示范版（136 词｜NCD 3 被动 + 复合介词 + 整句关系从句）}
Kære Hotel Kystblik

Jeg skriver denne klage \da{på baggrund af} et meget skuffende weekendophold på jeres hotel fra den 14. til den 16. november, hvor jeg havde bestilt værelse 204.

\da{Til trods for at} hotellet \dabold{beskrives} som roligt, opstod der flere alvorlige problemer under opholdet. For det første holdt høj musik fra gaden mig vågen begge nætter, fordi de gamle vinduer ikke lukkede tæt. Desuden var vandet i bruseren iskoldt lørdag morgen, og rengøringen \dabold{var ikke blevet udført} ordentligt.

Da jeg henvendte mig i receptionen, reagerede medarbejderen desværre uhøfligt. I stedet for at beklage svarede han surt, at støj var normalt i weekenden, \da{hvilket gjorde mig} meget skuffet.

Da jeg ikke fik den service, der \dabold{blev lovet} ved bestillingen, forventer jeg, at I refunderer halvtreds procent af beløbet inden fjorten dage.

Med venlig hilsen\\
Katrine Holm

\vspace{0.3em}
\rule{\linewidth}{0.4pt}\\[0.2em]
\small\textbf{\color{NCDNavy}[高分亮点 · 136 ord]}：\ding{172} 精准调用 NCD 3 书面复合介词 \textit{på baggrund af} 与让步结构 \textit{Til trods for at}；\ding{173} 灵活用出两种被动表达（现状描述用 \textit{-s} 被动 \textit{beskrives}，过去完成事件与动作分别用 \textit{var ikke blevet udført} 和 \textit{blev lovet}）；\ding{174} 末尾用整句关系从句 \textit{..., hvilket gjorde mig meget skuffet} 点评态度，措辞礼貌而坚定。
\end{emailbox}

\begin{emailbox}{Opgave 2 范文一：B1 稳过版（111 词｜超过 90 词底线 · 四问全答 · 自然亲切）}
Hej Astrid

Tak for din søde mail! Det er dejligt at høre fra dig igen.

Ja, jeg er lige begyndt på et aftenkursus i projektledelse hver tirsdag aften. Jeg valgte det, fordi jeg gerne vil lære nye ting og få mere ansvar på mit arbejde næste år.

Det kan godt være travlt at nå både fuldtidsarbejde, danskundervisning og fritid i hverdagen. Derfor laver jeg altid en ugeplan om søndagen og forbereder aftensmad til flere dage ad gangen.

Når det bliver weekend, vil jeg helst slappe af. Så slukker jeg arbejdsmailen, går en lang tur i skoven og mødes med venner til kaffe. Hvad laver du selv for tiden?

Mange hilsner\\
Rasmus

\vspace{0.3em}
\rule{\linewidth}{0.4pt}\\[0.2em]
\small\textbf{\color{NCDTeal}[考官点评 · 111 ord]}：稳超官方 90 词最低门槛，行文流畅自然。依次应答了 Astrid 全部问题（参加周二项目管理夜校、明年职场晋升动机、周日列周计划与提前备餐的时间管理法、周末关邮箱逛森林与喝咖啡放松），并在结尾反问互动，语气完全契合挚友通信语域。
\end{emailbox}

\begin{emailbox}{Opgave 2 范文二：B1+ 高分示范版（137 词｜条件倒装/方式状语 + 形式宾语 + 从句切换）}
Kære Astrid

Hvor er det dejligt at høre fra dig! Det passer, at jeg for nylig startede på et spændende aftenkursus i digital kommunikation. Jeg valgte holdet, fordi jeg ønsker at styrke mine faglige kompetencer på det danske arbejdsmarked, \da{selvom mine hverdage allerede er} ret fyldte.

Til tider kræver det god planlægning at få fuldtidsjob, sprogskole to aftener om ugen og fritid til at hænge sammen. Men \da{ved at skrive} alle aftaler i kalenderen og lægge mobilen væk efter klokken ni \dabold{når jeg} stadig både lektier, madlavning og søvn, selvom det \dabold{ikke altid er} nemt.

I weekenden lægger jeg derfor vægt \da{på, at jeg får koblet} helt af fra arbejdet. Jeg bager frisk brød lørdag morgen, laver rolig yoga derhjemme og bruger søndagen i naturen. Skal vi ikke snart mødes til en kop kaffe?

Kærlig hilsen\\
Rasmus

\vspace{0.3em}
\rule{\linewidth}{0.4pt}\\[0.2em]
\small\textbf{\color{NCDNavy}[高分亮点 · 137 ord]}：\ding{172} 方式状语短语 \textit{ved at skrive ... og lægge ...} 置句首直接触发主句 V2 倒装 \textit{når jeg}；\ding{173} 准确驾驭“介词 + 从句”结构（\textit{lægger vægt på, at...}）及让步从句中副词前置（\textit{selvom det ikke altid er nemt}）；\ding{174} 词汇丰富地道（\textit{faglige kompetencer, hænge sammen, koble af}），自然冲击 B1 上限高分。
\end{emailbox}

% =====================================================================
% 4.4 口语示范逐字稿、核心 IPA 与双人四图互动全真逐字稿
% =====================================================================
\subsubsection*{4.4 口语示范逐字稿、核心 IPA 与 Opgave 2 双人四图互动逐字稿}

\begin{transbox}
\small
\textbf{\color{NCDNavy}[Mindmap 2A 口语陈述范文：Mine sunde og usunde vaner｜150 词 · 约 1.5 分钟]}\\
\textit{Jeg vil gerne fortælle om mine sunde og usunde vaner \ipa{ˈvæːnɐ}, som spiller en stor rolle i min hverdag. På den sunde side starter jeg næsten altid dagen med havregryn \ipa{ˈhɑu̯ʁɐˌɡʁyːˀn} og frisk frugt, og til aftensmad spiser vi mange grøntsager \ipa{ˈɡʁœnˌsæːɐ}, fuldkornsbrød og fisk to gange om ugen. Jeg drikker også masser af vand i stedet for sodavand. Alligevel har jeg nogle usunde vaner, for når jeg er træt om aftenen, spiser jeg tit lidt for meget chokolade \ipa{ɕoɡoˈlæːðə} eller fredagsslik, og på kontoret drikker jeg nogle gange fire kopper kaffe. Med hensyn til motion \ipa{moˈɕoːˀn} cykler jeg otte kilometer til arbejde hver dag, svømmer \ipa{ˈsvømɐ} om onsdagen og går tur i skoven om søndagen. Jeg har aldrig røget, og jeg drikker kun et glas vin i weekenden. Til gengæld burde jeg forbedre min søvn \ipa{ˈsœu̯ˀn} ved at lægge mobilen væk klokken ti. Alt i alt er jeg tilfreds med min livsstil \ipa{ˈliusˌsdiːˀl}, fordi balancen fungerer godt.}
\end{transbox}

\paragraph*{Mindmap 2A Del 2 考官追问实战示范对话：}
\dialogturn{1}{Eksaminator}{Du siger, at du spiser lidt for meget chokolade om aftenen -- hvornår har du mest lyst til usunde ting, og hvorfor?}{chokolade \ipa{ɕoɡoˈlæːðə} · usunde \ipa{ˈuˌsʌnə}}{你说你晚上会吃有点多的巧克力——你什么时候最想吃不健康的东西，为什么？}
\dialogturn{2}{Kursist}{Det sker næsten altid efter klokken ni om aftenen, når jeg endelig sidder stille i sofaen efter en lang arbejdsdag. Jeg tror, det handler om, at jeg er træt i hovedet og bare vil belønne mig selv lidt.}{arbejdsdag \ipa{ˈɑːbɑi̯ˀdsˌdæːˀ} · belønne \ipa{beˈlœnˀə}}{几乎总是在晚上九点以后，经历漫长工作终于安静坐在沙发上的时候。我觉得这是因为脑子累了，想稍微犒劳一下自己。}
\dialogturn{3}{Eksaminator}{Hvorfor kan du godt lide at svømme og cykle til arbejde, og hvordan holder du motivationen om vinteren?}{motivationen \ipa{motˢivaˈɕoːˀnən}}{为什么你喜欢游泳和骑车上班，冬天又是怎么保持动力的？}
\dialogturn{4}{Kursist}{Når jeg cykler, vågner jeg godt op om morgenen, og i svømmehallen slapper hele kroppen af. Om vinteren har jeg varmt regntøj og lygter på cyklen, og fordi jeg har en fast aftale med en ven i svømmehallen hver onsdag, bliver jeg aldrig hjemme.}{svømmehallen \ipa{ˈsvøməˌhalˀən} · lygter \ipa{ˈløɡdɐ}}{骑车能让我早上彻底清醒，在游泳馆全身都能放松。冬天我备好了保暖雨衣和车灯，加上每周三和朋友约好了去游泳，所以我从不缺席。}
\dialogturn{5}{Eksaminator}{Du nævner, at du lægger telefonen væk klokken ti -- har det gjort en forskel for din søvn?}{forskel \ipa{ˈfʌːˌsɡɛlˀ} · søvn \ipa{ˈsœu̯ˀn}}{你提到你十点把手机放一边——这对你的睡眠带来了改变吗？}
\dialogturn{6}{Kursist}{Ja, en ret stor forskel faktisk. Før lå jeg og læste nyheder til midnat, så det var svært at falde i søvn. Nu læser jeg tyve minutter i en almindelig bog i stedet, og derfor sover jeg meget hurtigere.}{nyheder \ipa{ˈnyˌheðˀɐ} · almindelig \ipa{alˈmenˀəli}}{是的，其实改变还挺大的。以前我会躺着刷新闻到半夜，很难入睡。现在我改看二十分钟纸质书，因此入睡快得多。}
\dialogturn{7}{Eksaminator}{Hvad synes du generelt er det vigtigste for at leve et sundt liv i en travl hverdag i Danmark?}{generelt \ipa{ɡenəˈʁɑlˀd}}{总体而言，你认为在丹麦忙碌的日常生活中过上健康生活最重要的是什么？}
\dialogturn{8}{Kursist}{Jeg synes, det vigtigste er at have en god balance og nogle enkle rutiner, for eksempel at cykle til arbejde. Hvis reglerne er alt for strenge, giver man hurtigt op igen.}{balance \ipa{baˈlɑŋsə} · strenge \ipa{ˈsdʁɛŋə}}{我觉得最重要的是保持良好平衡并养成简单习惯，比如骑车上班。如果规矩订得太苛刻，人很快就会放弃。}

\vspace{0.3em}
\begin{transbox}
\small
\textbf{\color{NCDNavy}[Mindmap 2B 口语陈述范文：Mine digitale vaner｜146 词 · 约 1.5 分钟]}\\
\textit{Nu vil jeg gerne fortælle lidt om mine digitale vaner. I gennemsnit bruger jeg cirka to og en halv time på min telefon om dagen, men jeg har sat en skærmtidsgrænse \ipa{ˈsɡɛɐ̯mˌtˢiðsˌɡʁɛnsə} på flere apps, og ved aftensmaden ligger mobilen i tasken. På sociale medier \ipa{soˈɕæːlə ˈmeːðiɐ} bruger jeg mest LinkedIn til arbejde og kigger lidt på Instagram i bussen, men jeg har slået alle notifikationer \ipa{notˢifikaˈɕoːˀnɐ} fra for ikke at blive forstyrret. Jeg bruger også telefonen aktivt til at følge med i samfundet: Om morgenen læser jeg DR Nyheder, og når jeg cykler på arbejde, lytter jeg til danske podcasts for at øve udtalen \ipa{ˈuðˌtˢæːlən}. Derudover bruger jeg MitID næsten dagligt til netbank, e-Boks og borger.dk, mens jeg undgår impulskøb \ipa{emˈpʰulˀsˌkʰøːˀb} online. På arbejdet har vi ofte møder i Teams, hvilket sparer transporttid, men kræver skærmpauser. Generelt er jeg glad for digitale løsninger, selvom jeg gerne vil have flere skærmfrie aftener.}
\end{transbox}

\paragraph*{Mindmap 2B Del 2 考官追问实战示范对话：}
\dialogturn{1}{Eksaminator}{Du fortæller, at du har slået notifikationer fra -- hvordan har det påvirket din hverdag på arbejde?}{påvirket \ipa{ˈpʰʌˌviɐ̯ɡəð}}{你说你关掉了消息通知——这对你的日常工作有什么影响？}
\dialogturn{2}{Kursist}{Det har hjulpet rigtig meget på min koncentration. Før kiggede jeg hele tiden på skærmen, hver gang telefonen vibrerede, så det tog lang tid at blive færdig med en opgave. Nu tjekker jeg kun beskeder i min frokostpause.}{koncentration \ipa{kʰʌnsəntˢʁaˈɕoːˀn}}{这对我集中注意力帮助极大。以前手机一震动我就看屏幕，一项任务拖很久才做完；现在我只在午休时看消息。}
\dialogturn{3}{Eksaminator}{Hvad synes du om, at næsten alting i Danmark foregår digitalt med MitID og e-Boks?}{foregår \ipa{ˈfɒːˌɡʌˀ}}{几乎所有事在丹麦都要用 MitID 和 e-Boks 线上办理，你对此怎么看？}
\dialogturn{4}{Kursist}{For mig er det super praktisk, fordi jeg kan ordne lægetid, bank og breve fra kommunen hjemmefra på fem minutter. Men til trods for at det er smart, synes jeg også, det kan være svært for ældre mennesker eller nyankomne udlændinge.}{praktisk \ipa{ˈpʰʁɑɡtˢisɡ} · udlændinge \ipa{ˈuðˌlɛneŋə}}{对我来说非常方便，在家五分钟就能搞定预约医生、银行和市政信件。不过尽管很智能，我也觉得这对老年人或新来的外国人可能有点难。}
\dialogturn{5}{Eksaminator}{Vil du fortælle lidt mere om, hvordan du bruger danske podcasts eller apps til at lære dansk?}{udtalen \ipa{ˈuðˌtˢæːlən}}{能再多讲讲你是怎么用丹麦语播客或软件来学丹麦语的吗？}
\dialogturn{6}{Kursist}{Ja, selvfølgelig! Jeg lytter til en kort nyhedspodcast på femten minutter hver morgen, hvor de taler tydeligt. Hvis jeg hører et nyt ord, skriver jeg det ind i min ordbogsapp bagefter, så jeg kan øve det igen om aftenen.}{nyhedspodcast \ipa{ˈnyheðsˌpʰʌdˌkʰɑsd}}{当然！我每天早上听十五分钟语速清晰的短新闻播客。如果听到新词，之后就会记进词典应用里晚上再复习。}
\dialogturn{7}{Eksaminator}{Tror du, vi bruger telefonen for meget i dag, eller er der mest fordele ved digitale løsninger?}{fordele \ipa{ˈfɒːˌdeːlə}}{你觉得如今人们手机用得太多了吗，还是说数字化方案利大于弊？}
\dialogturn{8}{Kursist}{Jeg tror faktisk, mange af os bruger den lidt for meget om aftenen, fordi man let glemmer tiden med sociale medier. Så selvom teknologien har store fordele på arbejde og i samfundet, er det sundt selv at sætte klare grænser.}{grænser \ipa{ˈɡʁɛnsɐ}}{我确实觉得很多人晚上用得过多，因为刷社交媒体容易忘了时间。所以尽管科技在工作和社会中优势巨大，自己设好明确界限才健康。}

\paragraph*{Mundtlig Opgave 2（EMNE: Fritid og motion）Del 1 双人四图互动全真逐字稿：}
\dialogturn{1}{Kursist A}{Nå, vores emne er 'Fritid og motion'. Hvis vi kigger på de fire billeder her, hvad synes du så bedst om? Kan du godt lide at træne i et fitnesscenter, som på det første billede?}{fitnesscenter \ipa{ˈfidnəsˌsɛnˀdɐ}}{好，我们的话题是“休闲与运动”。看看这四张图，你最喜欢哪一种？你喜欢像第一张图那样在健身房训练吗？}
\dialogturn{2}{Kursist B}{Altså, jeg går faktisk i fitnesscenter to gange om ugen, fordi det er meget fleksibelt. Jeg kan træne efter arbejde, også når det regner eller er mørkt udenfor. Men nogle gange bliver det lidt kedeligt at stå på et løbebånd. Hvad med dig -- foretrækker du fitnesscenter eller at løbe ude i skoven som på billede to?}{løbebånd \ipa{ˈløːbəˌbʌnˀ} · skoven \ipa{ˈsɡʌu̯ˀən}}{其实我每周去两次健身房，因为时间非常灵活。哪怕外面下雨或天黑，我下班后也能练。不过有时在跑步机上跑也有点无聊。你呢——你更喜欢健身房，还是像图二那样在森林里户外跑？}
\dialogturn{3}{Kursist A}{Jeg kan helt klart bedst lide at løbe ude i skoven! Selvom man selvfølgelig bliver lidt våd i dårligt vejr, giver den friske luft og de grønne træer mig meget mere energi end et fitnesscenter med høj musik. Og så er det helt gratis. Hvad tænker du om holdsport i en hal, som vi ser på det tredje billede?}{friske luft \ipa{ˈfʁɛsɡə lufd} · holdsport \ipa{ˈhʌlˀˌsbʌːd}}{我明显更喜欢在森林里跑步！虽然坏天气难免淋湿，但新鲜空气和绿树比放着响亮音乐的健身房更能给我充电，而且还免费。你怎么看图三那种在体育馆里进行的团体球类运动？}
\dialogturn{4}{Kursist B}{Jeg elsker holdsport! Jeg spiller badminton hver torsdag i en lokal klub. Det bedste ved holdsport er det sociale fællesskab, for man snakker, griner og glemmer næsten, at man dyrker hård motion. Ulempen er bare, at man skal passe faste tider hver uge. Kunne du selv tænke dig at spille badminton eller volleyball?}{fællesskab \ipa{ˈfɛləsˌsɡæːˀb} · ulempen \ipa{ˈuˌlɛmbən}}{我超爱团体运动！我每周四都在当地俱乐部打羽毛球。团体运动最棒的是社交氛围，大家边说边笑，几乎忘了自己正在剧烈运动。唯一缺点是要迁就每周固定时间。你自己愿意试试打羽毛球或排球吗？}
\dialogturn{5}{Kursist A}{Ja, det lyder faktisk rigtig hyggeligt, især hvis man kan melde sig på et begynderhold sammen med en ven. Men når jeg har haft en meget stressende uge, har jeg nogle gange bare brug for ro. Derfor kan jeg også rigtig godt lide billede fire med yoga og udstrækning derhjemme. Laver du nogensinde rolige øvelser i stuen?}{stressende \ipa{ˈsdʁɛsənə} · udstrækning \ipa{ˈuðˌsdʁɛɡneŋ}}{是的，听起来真的挺惬意，尤其是和朋友一起报个初级班。不过经历高压的一周后，有时我只想静一静，因此我也很喜欢图四那样在家做瑜伽和拉伸。你在客厅做过舒缓训练吗？}
\dialogturn{6}{Kursist B}{Ja, engang imellem laver jeg strækøvelser efter arbejde, når min ryg er stiv. Det er rart, at man ikke behøver at tage nogen steder hen og kun bruger ti minutter på en måtte. Til gengæld er det lidt svært at holde koncentrationen alene derhjemme, så jeg foretrækker stadig holdsport eller fitness.}{måtte \ipa{ˈmʌdə} · stiv \ipa{ˈsdiːu̯ˀ}}{会的，下班后如果背部僵硬，我偶尔会做拉伸。不用出门、垫子上十分钟就能练完很方便。不过一个人在家不太容易保持专注，所以我还是更偏爱团体运动或健身房。}

\paragraph*{Mundtlig Opgave 2 Del 2（个人追问 3 题高分示范）：}
\dialogturn{Q1}{Tester}{Hvordan kan du bedst lide at slappe af eller dyrke motion i din fritid? Hvorfor?}{slappe af \ipa{ˈslɑbə ɑˀ}}{你空闲时最喜欢如何放松或锻炼？为什么？}
\dialogturn{A1}{Kursist}{Jeg kan allerbedst lide at løbe en lang tur i skoven alene efter arbejde, fordi det hjælper mig med at tømme hovedet for arbejdsstress. Når man hører fuglene og ser årstiderne skifte, føles det slet ikke som hård træning, men mere som en pause for hjernen.}{årstiderne \ipa{ˈɒːsˌtˢiðˀɐnə}}{我最喜欢下班后独自去森林跑长距离，这能帮我清空工作压力。听着鸟叫、看着四季变化时，完全不像苦练，更像给大脑放了个假。}
\dialogturn{Q2}{Tester}{Hvornår og hvor ofte har du tid til fritidsinteresser i en almindelig uge?}{almindelig \ipa{alˈmenˀəli}}{在平常的一周里，你一般什么时候、多久有时间投入业余爱好？}
\dialogturn{A2}{Kursist}{I en normal uge har jeg tid til motion omkring tre gange. Jeg prøver at løbe en kort tur tirsdag og torsdag aften, når aftensmaden er spist, og så går jeg en længere tur eller laver yoga søndag formiddag, hvor der ikke er så mange faste aftaler.}{formiddag \ipa{ˈfʌːmiˌdæːˀ}}{正常一周里我大概有三次锻炼时间。周二和周四吃完晚饭后会试着短跑一次，周日上午没什么固定安排，就会走远一点或做做瑜伽。}
\dialogturn{Q3}{Tester}{Kan du fortælle lidt om en fritidsaktivitet eller sport, du har prøvet eller gerne vil prøve?}{aktivitet \ipa{ɑɡtˢiviˈtˢeːˀd}}{你能讲讲你尝试过或想尝试的一项休闲活动或运动吗？}
\dialogturn{A3}{Kursist}{Jeg meldte mig til et begynderhold i badminton for et halvt år siden med min kollega Morten. I starten var det svært at ramme bolden i det høje tempo, men vores træner var meget tålmodig, og i dag elsker jeg det, fordi man både får sved på panden og nye venner.}{tålmodig \ipa{tˢʌlˈmoːˀði}}{半年前我和同事莫滕一起报了羽毛球零基础班。一开始节奏很快、很难击中球，但教练很有耐心，现在我非常喜欢这项运动，既能出汗又能交到新朋友。}
